%%%%%%%%%%%%%%%%%%%%%%%%%%%%%%%%%%%%%%%%%%%%%%%%%%%%%%%%%%%%%%%%%%%%%%%%%%%%%%%%
%\section{Finite $P_z$ Corrections for Unpolarized PDFs}\label{sec:corrections}
%%%%%%%%%%%%%%%%%%%%%%%%%%%%%%%%%%%%%%%%%%%%%%%%%%%%%%%%%%%%%%%%%%%%%%%%%%%%%%%%

In this section, we detail the procedure to implement the $P_z$ corrections
needed to extract the physical $q(x,\mu)$ from the quasi-distribution
$\tilde{q}(x,\Lambda,P_z)$ computed from the lattice. We first explain each of the
corrections and then summarize the procedure at the end of the section.

%%%%%%%%%%%%%%%%%%%%%%%%%%%%%%%%%%%%%%%%%%%%%%%%%%%%%%%%%%%%%%%%%%%%%%%%%%%%%%%%
\subsection{One-Loop Matching}

In the limit $P_z \rightarrow \infty$, the matching becomes the most
important $P_z$ correction. The factor
$Z\left(\xi=\frac{x}{y},\frac{\mu}{P_z},\frac{\Lambda}{P_z}\right)$
has been computed up to one loop in
Ref.~\cite{Xiong:2013bka} using a momentum-cutoff regulator instead of a
lattice regulator. Therefore, this $Z$ factor is accurate up to the
leading logarithm but not for the numerical constant. To determine this
constant, a lattice perturbation theory calculation using the same lattice
action is required.

At tree level, the $Z$ factor is just a delta function,
%
\begin{equation}
Z(\xi) = \delta (\xi-1) + \frac{\alpha_s}{2\pi} \overline{Z}(\xi)
  + \mathcal{O}\left(\alpha_s^2 \right),
\end{equation}
%
such that
%
\begin{equation}
\widetilde{q}(x) \simeq q(x)
  + \frac{\alpha_s}{2\pi} \int \frac{dy}{\left\vert y\right\vert}\,
    \overline{Z}\!\left(\frac{x}{y}\right)q(y).
\end{equation}
%
Since the difference between $\widetilde{q}(x)$ and $q(x)$ starts at
the loop level, we can rewrite the above equation as
%
\begin{equation}
q(x) \simeq \widetilde{q}(x)
  - \frac{\alpha_s}{2\pi} \int \frac{dy}{\left\vert y\right\vert}\,
    \overline{Z}\!\left(\frac{x}{y}\right)\widetilde{q}(y)
\label{q}
\end{equation}
%
with an error of $\mathcal{O}\left(\alpha_s^2\right)$~\cite{Ma:2014jla}.

$\overline{Z}(\xi)$ is a singular function of $\xi$ with terms like
$a/\left(1-\xi\right)^2$ and
$\left(b\ln \left\vert 1-\xi \right\vert + c\right)/(1-\xi)$~\cite{Xiong:2013bka}. The value of $\left\vert \xi \right\vert$ is not bounded by unity when $P_z$ is
finite. This is because $\overline{Z}$ describes the evolution of partons, which
includes the gluon emission process where a mother quark splits into a daughter
quark and a gluon. When $P_z$ is finite, the gluon can travel in the opposite
direction from the mother quark, which makes the momentum fraction of the
daughter quark bigger than that of the mother quark; hence, $\xi$ can be bigger
than one. To show that these singular terms are not harmful, we use the fact
that $\overline{Z}$ has the structure
%
\begin{equation}
\overline{Z}(\xi) = \left(Z^{(1)}(\xi) - C\delta(\xi-1)\right),
\end{equation}
%
with the first term coming from gluon emission and the second term from the quark
self-energy diagram. $C=\int_{-\infty}^{\infty}\!d\xi^{\prime}\,Z^{(1)}(\xi
^{\prime})$ such that $\int\!d\xi\,\overline{Z}(\xi)=0$ and particle-number
conservation is satisfied. Using this, Eq.~\ref{q} becomes
%
\begin{equation}
q(x) \simeq \widetilde{q}(x)
  - \frac{\alpha_s}{2\pi} \int_{-\infty}^{\infty}\!dy\,
  \left[ Z^{(1)}\!\left(\frac{x}{y}\right)
    \frac{\widetilde{q}(y)}{\left\vert y\right\vert}
  - Z^{(1)}\!\left(\frac{y}{x}\right)
    \frac{\widetilde{q}(x)}{\left\vert x\right\vert}\right] .  \label{cov}
\end{equation}
%
Then, when $\xi=x/y$ is close to one, the double poles
$1/\left(1-\xi\right)^2$ in $Z^{(1)}(\xi)$ cancel out. The single pole
$\left(b\ln \left\vert 1-\xi \right\vert + c\right)/(1-\xi)$ is odd in $(y-x)$,
which is not an endpoint singularity in $y$ because $\xi$ is not bounded
by unity. Therefore, the integral and $q(x)$ are finite.

Another observation is that the double-pole structure in $Z^{(1)}$ makes
$\widetilde{q}(x)$ behave like $\alpha_s/x^2$ at large $x$. This can be
seen by approximating $\widetilde{q}(y)$ in the integral by $q(y)$ with $y$
bounded by unity. So even though $\widetilde{q}(x)$ is finite, its higher
moments are not. This is consistent with the fact that moments of
$\tilde{q}(x)$, $q(x)$ and $Z(x)$ satisfy
%
\begin{equation}\label{momentrelation}
\left\langle x^n \right\rangle_{\tilde{q}} =
  \left\langle x^n \right\rangle_Z \left\langle x^n \right\rangle_q,
\end{equation}
%
where
\begin{equation}
\langle x^{n}\rangle_{f}=\int_{-\infty}^{\infty}dx\,x^{n}f(x).
\end{equation}
Note that for $f=q$ this is identical to the usual moments since $q(x)$ vanishes outside $[-1,1]$. Eq.~\ref{momentrelation} can be obtained by taking moments of Eq.~\ref{Z}, and the divergence
of higher moments of $Z$ implies the divergences of higher moments of $\tilde{q}$.

%%%%%%%%%%%%%%%%%%%%%%%%%%%%%%%%%%%%%%%%%%%%%%%%%%%%%%%%%%%%%%%%%%%%%%%%%%%%%%%%
\subsection{$M^{2n}/P_z^{2n}$ Correction}

In this subsection, we show how the $M^2/P_z^2$ correction to all
orders (denoted as $M^{2n}/P_z^{2n}$) can be computed exactly. We will
make use of the property
%
\begin{equation}
\lambda_{\mu_1}\cdots \lambda_{\mu_n} P^{(\mu_1}\cdots P^{\mu_n)} =
  \lambda_{(\mu_1}\cdots \lambda_{\mu_n)} P^{\mu_1}\cdots P^{\mu_n},
\end{equation}
%
where $(...)$ means the indices enclosed are symmetric and traceless.

A useful identity is
%
\begin{equation}
\lambda_{(\mu_1}\cdots \lambda_{\mu_n)} =
  \sum_{i=0}^{i_\text{max}} B_{n,i} \left(\lambda^2\right)^i
    \left(\frac{\partial^2}{\partial\lambda_\alpha \partial\lambda^\alpha}\right)^i
    \lambda_{\mu_1}\cdots \lambda_{\mu_n},  \label{a1}
\end{equation}
%
where $i_\text{max}=\frac{n-\text{Mod}[n,2]}{2}$ and $B_{n,0}=1$.
The $B$ coefficients can be determined by the tracelessness of
$\lambda_{(\mu_1}\cdots \lambda_{\mu_n)}$ which implies
%
\begin{equation}
g^{\mu_1 \mu_2}P^{\mu_3}\cdots P^{\mu_n}
  \lambda_{(\mu_1}\cdots \lambda_{\mu_n)} = 0,
\end{equation}
%
or
%
\begin{equation}
\sum_{i=0}^{i_\text{max}} B_{n,i} \left(\lambda^2\right)^i
  \left(\frac{\partial^2}{\partial\lambda_\alpha \partial\lambda^\alpha}\right)^i
  \lambda^2 \left(\lambda \cdot P\right)^{n-2} = 0.  \label{C}
\end{equation}
%
The left-hand side of this equation is a polynomial of powers of $\lambda^2$
with each term involving at most two $B$ coefficients. Then, the identity of
Eq.~\ref{C} yields the following recurrence relation:
%
\begin{equation}
B_{n,i} = -\frac{B_{n,i-1}}{4i(n-i+1)}.  \label{a3}
\end{equation}

To implement the $M^{2n}/P_z^{2n}$ correction, we first compute the ratio of the moments 
%
\begin{align}
K_n  &\equiv\frac{\left\langle x^{n-1}\right\rangle_{\tilde q}}{\left\langle x^{n-1}\right\rangle_q}=\frac{\lambda_{(\mu_1}\cdots \lambda_{\mu_n)}
  P^{\mu_1}\cdots P^{\mu_n}}{\lambda_{\mu_1}\cdots \lambda_{\mu_n}
  P^{\mu_1}\cdots P^{\mu_n}}  \notag \\
&= \sum_{i=0}^{i_\text{max}} C_{n-i}^i c^i   \label{Kn},
\end{align}
%
where $C$ is the binomial function and
$c=-\lambda^2 M^2/4\left(\lambda\cdot P\right)^2 = M^2/4 P_z^2$
with $\lambda^\mu = (0,0,0,-1)$ and $\lambda \cdot P=P_z$. 

As shown in Appendix B, the above factors can be converted to the following relation between PDFs
\begin{align}
q(x)&=\sqrt{1+4c}\sum_{n=0}^\infty \frac{f_-^n}{f_+^{n+1}}\Big[(1+(-1)^n)\tilde q\Big(\frac{f_+^{n+1}x}{2f_-^n}\Big)+(1-(-1)^n)\tilde q\Big(\frac{-f_+^{n+1}x}{2f_-^n}\Big)\Big]\non\\
&=\sqrt{1+4c}\sum_{n=0}^\infty \frac{(4c)^n}{f_+^{2n+1}}\Big[(1+(-1)^n)\tilde q\Big(\frac{f_+^{2n+1}x}{2(4c)^n}\Big)+(1-(-1)^n)\tilde q\Big(\frac{-f_+^{2n+1}x}{2(4c)^n}\Big)\Big],
\end{align}
where $f_{\pm}=\sqrt{1+4c}\pm 1$. Unlike the mass-correction expression of Ref.~\cite{Alexandrou:2015rja}, particle number is conserved in this expression.

%%%%%%%%%%%%%%%%%%%%%%%%%%%%%%%%%%%%%%%%%%%%%%%%%%%%%%%%%%%%%%%%%%%%%%%%%%%%%%%%
\subsection{$\Lambda_\text{QCD}^2/P_z^2$ Correction}

This correction comes from the trace part on the left-hand side of Eq.~\ref{O3},
which is a twist-4 effect and can be implemented by adding a
$\tilde{q}_\text{twist-4}$ contribution to $\tilde{q}$, such that
%
\begin{equation}
\tilde{q}(x,\Lambda,P_z) \rightarrow
  \tilde{q}(x,\Lambda,P_z) + \tilde{q}_\text{twist-4}(x,\Lambda,P_z).
\end{equation}
%
As derived in Appendix~C,
%
\begin{equation}
\tilde{q}_\text{twist-4}(x,\Lambda,P_z) =
  \frac{1}{8\pi}\int_{-\infty}^\infty\!dz\,\Gamma_0 \left(-ixzP_z\right)
  \left\langle P\left\vert \mathcal{O}_\text{tr}(z)\right\vert P\right\rangle,
  \label{t4-1}
\end{equation}
%
where $\Gamma_0$ is the incomplete Gamma function and
%
\begin{align}
\mathcal{O}_\text{tr}(z) &=
  \int_0^z\!dz_1\,\bar{\psi}(0) \Big[ \gamma^\nu \Gamma\left(0,z_1\right)
    D_\nu\Gamma \left(z_1,z\right) \notag \\
&{} + \int_0^{z_1}\!dz_2\, \lambda \cdot \gamma
  \Gamma\left(0,z_2 \right) D^\nu \Gamma \left(z_2,z_1\right)
  D_\nu \Gamma\left(z_1,z\right) \Big] \psi (z\lambda).  \label{t4-2}
\end{align}
%
Instead of computing these corrections directly on the lattice, we only
parametrize and fit them as a $1/P_z^2$ correction after we have
removed other leading-$P_z$ corrections.

%%%%%%%%%%%%%%%%%%%%%%%%%%%%%%%%%%%%%%%%%%%%%%%%%%%%%%%%%%%%%%%%%%%%%%%%%%%%%%%%
\subsection{Summary of the Finite $P_z$ Corrections}

Here we summarize the procedure needed to implement finite-$P_z$ corrections,
and use the unpolarized $u(x)-d(x)$ PDF as an example. We focus on the flavor-nonsinglet PDF such that there is no mixing with the gluon PDF.
The generalization to the polarized case will be shown in the next section.

We start with the computation of the equal-time correlator on a Euclidean
lattice:
%
\begin{equation}
h_\text{lat}\left(z,P_z,\Lambda\right) =
  \frac{1}{2P_z} \left\langle P\left\vert \bar{\psi}(0) \gamma_z
  \left(\prod_n U_z (n\hat{z})\right) \psi(z) \right\vert P\right\rangle,
  \label{hlat}
\end{equation}
%
which is the lattice version of Eq.~\ref{eq:h-definition} with
$\lambda^\mu=(0,0,0,-1)$ and $U_\mu$ is a discrete gauge link in the $\mu$
direction. $h_\text{lat}\left(0\right)$ is the total quark number of
the hadronic state. The quasi-PDF is the Fourier transform of $h_\text{lat}$:
%
\begin{equation}
\widetilde{q}(x,P_z,\Lambda) \equiv
  \int_{-\infty}^\infty \frac{dz\,P_z}{2\pi} e^{ixzP_z} h_\text{lat}.
\end{equation}

The second step is to implement the one-loop matching to convert
$\widetilde{q}(x,P_z,\Lambda)$ in the lattice scheme to $q_I(x,P_z,\mu)$
in the $\overline{\text{MS}}$ scheme:

\begin{equation}
q_I(x) \simeq \widetilde{q}(x) -
  \frac{\alpha_s}{2\pi} \int_{-\infty}^\infty\!dy\,
  \left[ Z^{(1)}\!\left(\frac{x}{y}\right)
    \frac{\widetilde{q}(y)}{\left\vert y\right\vert} - 
    Z^{(1)}\!\left(\frac{y}{x}\right)
    \frac{\widetilde{q}(x)}{\left\vert x\right\vert} \right] .  \label{Z(1)}
\end{equation}
%
We will use the $Z^{(1)}$ factor derived in Ref.~\cite{Xiong:2013bka} (also listed in Appendix~A for completeness),
which matches the momentum cutoff scheme with the $\overline{\text{MS}}$
scheme. This $Z^{(1)}$ factor is accurate up to the leading logarithm
but not for the numerical constant. One should replace the momentum cutoff
calculation with a lattice perturbation theory calculation to get the
correct numerical constant in the future.

The third step is to remove the $\mathcal{O}(M^{2n}/P_z^{2n})$
correction from $q_I (x,P_z,\mu)$:
%
\begin{align}
q_{II}(x)&=\sqrt{1+4c}\sum_{n=0}^\infty \frac{f_-^n}{f_+^{n+1}}\Big[(1+(-1)^n) q_I\Big(\frac{f_+^{n+1}x}{2f_-^n}\Big)+(1-(-1)^n)q_{I}\Big(\frac{-f_+^{n+1}x}{2f_-^n}\Big)\Big]\non\\
&=\sqrt{1+4c}\sum_{n=0}^\infty \frac{(4c)^n}{f_+^{2n+1}}\Big[(1+(-1)^n) q_I\Big(\frac{f_+^{2n+1}x}{2(4c)^n}\Big)+(1-(-1)^n) q_I\Big(\frac{-f_+^{2n+1}x}{2(4c)^n}\Big)\Big].
\end{align}
%
This gives rise to $q_{II}(x,P_z,\mu)$ whose remaining
$\mathcal{O}\left(\Lambda_\text{QCD}^2/P_z^2\right)$ correction can be removed by
adding the twist-4 contribution $\tilde{q}_\text{twist-4}$ defined in
Eqs.~\ref{t4-1}--\ref{t4-2}
%
\begin{equation}
q(x,\mu) = q_{II}(x,P_z,\mu) + \tilde{q}_\text{twist-4}(x,P_z,\mu).
\label{final step}
\end{equation}
%
$\tilde{q}_\text{twist-4}(x,P_z,\Lambda)$ can be computed on the lattice and in
principle another matching to the $\overline{\text{MS}}$ scheme is
required to obtain $\tilde{q}_\text{twist-4}(x,P_z,\mu)$. However, the difference
is $\mathcal{O}\left(\alpha_s\,\tilde{q}_\text{twist-4}\right)$ and hence negligible. In this work, the effect of $\tilde{q}_\text{twist-4}$ will be
parametrized, fitted to data and extrapolated $P_z \to \infty$ to obtain the $P_z$-independent left-hand side of
Eq.~\ref{final step}.
